\clearpage
\fancyhead[L]{\small D'AGUILA.}
\hypertarget{native-546-middle}{}\label{native-546-middle}
\begin{center}\small Göttingische gelehrte Anzeigen. Stück 102. S. 1009..1010. 27. Juni 1812.\\[3pt]\rule{.72\linewidth}{.3pt}\end{center}

\emph{Decouverte de l'orbite de la terre du point central de l'orbite du soleil, leur situation et leur forme; de la section du zodiaque, par le plan de l'équateur, et du mouvement concordant des deux globes. Avec figures. Par Mr.} \textsc{C. J. E. H. d'Aguila}, \emph{ancien élève du génie.} 1806. 432 Seiten in Oktav. Paris. \emph{De l'imprimerie de Délance.}

Auf Wunsch holen wir hier eine kurze Besprechung dieses Buches nach. Es ist zwar schon vor 6 Jahren erschienen, uns aber erst seit kurzem bekannt. Wir haben hier einen neuen Antikopernikaner, der die Sonne jährlich einen Kreis um die Erde beschreiben und die Planeten sich um die Sonne bewegen lässt (obwohl von diesen nur beiläufig und mit wenigen Worten die Rede ist). Der Erde schreibt er dagegen außer ihrer Rotation noch eine jährliche gleichförmige Bewegung in einem exzentrischen Kreis in der Ebene des Äquators zu. Den Radius dieses Kreises setzt er auf $\frac{1}{59}$ des Radius der jährlichen Sonnenbahn an. Diese wichtige Entdeckung, die der Verfasser selbst unbegreiflich nennt und einer unmittelbaren Inspiration zuschreibt, weshalb er sie auch Gott selbst widmet, soll das neunzehnte Jahrhundert verherrlichen und die Sekten von \textsc{Copernicus}, \textsc{Kepler} und \textsc{Newton}, auf die der Verfasser mit tiefer Verachtung herabsieht, vernichten. Das mag genügen, um das Buch neugierigen Lesern zu empfehlen, wenn auch nur als psychologische Merkwürdigkeit.
